\subsection{Content of an element of a free module}

Let A be a principal ideal domain, let L be a free A-module, and let x be an element of L. As f runs through the set \(L^{*}\) of linear forms on L, the elements \(f(x)\) form an ideal \(\mathfrak{c}_{\mathrm{L}}(x)\) of A, called the content of x in L. An element c of A is called a content of x in L if it generates the ideal \(\mathfrak{c}_{\mathrm{L}}(x)\) ; this amounts to saying that there exists a linear form f on L such that \(f(x)=c\) and that c divides \(g(x)\) for every linear form g on L. Let \((e_{i})_{i\in I}\) be a basis of L; put \(x=\sum a_{i}e_{i}, a_{i}\in A\) ; then the ideal \(\mathfrak{c}_{\mathrm{L}}(x)\) consists of sums \(\sum a_{i}b_{i}\) , as \((b_{i})\) runs through the set \(A^{I}\) ; it follows immediately that an element c of A is a content of x in L if and only if it is a gcd of the family \((a_{i})\) of coordinates of x.

We say that x is indivisible if \(\mathfrak{c}_{L}(x)=\mathbf{A}\) , that is if the coordinates of x with respect to a basis of L are setwise coprime.

Lemma 1. --- Let L be a free module over a principal ideal domain A and let x be an element of L. Then the following conditions are equivalent :

\begin{enumerate}
\def\labelenumi{(\roman{enumi})}
\item
  \(x\) is indivisible;
\item
  there exists a linear form \(f\) on \(L\) such that \(f(x) = 1\) ;
\item
  \(x\) is nonzero and the submodule \(\mathbf{A}x\) is a direct factor of \(\mathbf{L}\) ;
\item
  \(x\) is an element of some basis of \(\mathbf{L}\) .
\item
  \(\Rightarrow\) (ii): this follows from the definition.
\item
  \(\Rightarrow\) (iii): let \(f\) be a linear form on \(\mathbf{L}\) such that \(f(x) = 1\) ; then \(x \neq 0\) and the map \(y \mapsto f(y)x\) is a projector of \(\mathbf{L}\) , with image \(\mathbf{A}x\) .
\item
  \(\Rightarrow\) (iv): let \(L'\) be a complement of \(Ax\) in \(L\) , and let \(B'\) be a basis of \(L'\) (VII, p.~15, Cor. 2); then \(B' \cup \{x\}\) is a basis of \(L\) .
\item
  \(\Rightarrow\) (i): trivial.
\end{enumerate}

Remarks. --- 1) If \(x\) is a nonzero element of \(L\) and \(c\) a content of \(x\) , then there exists a unique element \(y\) of \(L\) such that \(x = cy\) ; denote this element \(x / c\) ; then \(x / c\) is an indivisible element of \(L\) .

\begin{enumerate}
\def\labelenumi{\arabic{enumi})}
\setcounter{enumi}{1}
\tightlist
\item
  The content \(\mathfrak{c}_{\mathrm{L}}(x)\) is the annihilator of the torsion module of \(\mathrm{L} / \mathrm{A}x\) .
\end{enumerate}

Let L be a free module over a principal ideal domain A and let M be a submodule of L; by VII, p.~4, Lemma 1, the family \(\mathfrak{c}_{\mathrm{L}}(x)\) , \(x \in M\) , has a maximal element; if \(M \neq \{0\}\) then such a maximal element is nonzero.

PROPOSITION 3. --- Let L be a free module over a principal ideal domain A and let M be a nonzero submodule of L. Let x be an element of M such that \(\mathfrak{c}_{\mathrm{L}}(x)\) is maximal among the contents of elements of M, let c be a content of x in L, and let f be a linear form on L such that \(f(x) = c\) .

\begin{enumerate}
\def\labelenumi{\alph{enumi})}
\item
  L is the direct sum of A(x/c) and the kernel K of f.
\item
  \(\mathbf{M}\) is the direct sum of \(\mathbf{A}\mathbf{x}\) and \(\mathbf{K} \cap \mathbf{M}\) .
\item
  \(g(M) \subset Ac\) for every linear form \(g\) on \(L\) .
\end{enumerate}

Put \(y = x / c\) ; clearly \(\mathbf{A}y \cap \mathbf{K} = \{0\}\) , since \(f(y) = 1\) . Furthermore, for each \(u \in L\) we have

\[
u = f (u) y + (u - f (u) y),
\]

with \(f(u) y \in \mathbf{A}y\) and \(u - f(u)y \in \mathbf{K}\) ; this proves \(a\) ). Now note that for \(u \in \mathbf{M}\) we have \(f(u) \in \mathbf{A}c\) : indeed, let \(u \in \mathbf{M}\) , and let \(d\) be a gcd of \(f(u)\) and \(c\) ; then there exist \(\lambda, \mu \in \mathbf{A}\) such that \(d = \lambda f(u) + \mu c = f(\lambda u + \mu x)\) ; hence the content of the element \(\lambda u + \mu x\) of \(\mathbf{M}\) divides \(d\) ; by the maximality of \(c\) , this implies that \(d\) is an associate of \(c\) , so \(f(u) \in \mathbf{A}c\) . Hence for all \(u\) in \(\mathbf{M}\) , we can write

\[
u = (f (u) / c) x + (u - (f (u) / c) x) \in \mathrm{A} x + (\mathrm{K} \cap \mathrm{M}),
\]

which shows \(b\) ). Finally, let \(g\) be a linear form on \(L\) ; by \(a\) ) there exists a scalar \(\alpha \in A\) and a linear form \(h\) on \(K\) such that \(g(u) = \alpha f(u) + h(u - f(u)y)\) ; thus by b) we have \(g(\mathbf{M}) \subset \mathbf{A}c + h(\mathbf{K} \cap \mathbf{M})\) . To prove c), it is therefore sufficient to prove that \(h(\mathbf{K} \cap \mathbf{M}) \subset \mathbf{A}c\) for every linear form h on K, or equivalently for every linear form h on L such that \(h(x) = 0\) ; now, if \(u \in K \cap M\) and d is a gcd of \(h(u)\) and c, then there exist \(\lambda\) , \(\mu \in A\) with \(d = \lambda h(u) + \mu c\) ; then \((f + h)(\lambda u + \mu x) = d\) , which implies as above that \(h(u) \in Ac\) , and c) follows.

